\chapter{Patterns: empty states}
\label{ch:empty}

\begin{aside}
What does your app show when the list is empty, the search
returns nothing, the user just started? Empty states deserve
design attention.
\end{aside}

\section{Three types}

\begin{itemize}[leftmargin=*]
  \item First use: ``Welcome! Try adding a thing.''
  \item Empty result: ``No matches for X. Try a different
    query.''
  \item Cleared: ``All done. Nothing to do here.''
\end{itemize}

Each calls for different copy.

\section{Visual}

A centred message. Maybe an ASCII illustration. A primary
action button.

\section{Don't show emptiness}

If the area is too small for a message, just leave it
blank with a hint elsewhere. Don't waste cells screaming
``empty''.

\section{Encouraging action}

The empty state often suggests what to do next: ``Click
here to add'', ``Press / to search''. The user is not
stuck.

\section{Recap}

\begin{recap}
\begin{itemize}[leftmargin=*]
  \item Different copy for first-use, empty result,
    cleared.
  \item Suggest the next action.
  \item Don't shout emptiness; small visuals.
\end{itemize}
\end{recap}

\section*{Workshop}

\begin{enumerate}[leftmargin=*]
  \item Identify every list in your app; design its
    empty state.
  \item Add empty-state messages.
\end{enumerate}
